\begin{textsample}{POS Dim 3 – global\_north\_2024\_02 – Score 10.00 – global\_north\_2024\_02/105761992.txt}  \label{ex:f3_pos_258}
Cook County prosecutors on Wednesday dropped charges against two Northwestern \textbf{students} amid controversy over the decision to pursue a criminal case against them for allegedly circulating a fake page of the \textbf{student} newspaper to \textbf{protest} the \textbf{school} stance on the crisis in Gaza and Israel.
The charging decision, as well as the newspaper ' s parent company ' s pursuit of a complaint, attracted harsh criticism from \textbf{students}, professors and community members, who blasted the move as an example of over-policing of Black \textbf{students} and an effort to silence pro-Palestinian voices that disproportionately affect people of color.
In a statement, Northwestern \textbf{University} said \textbf{campus} \textbf{police} acted on a complaint filed by \textbf{Students} Publishing company and investigated the distribution of fake Daily Northwestern news pages."The \textbf{University} supports the Cook County State ' s \textbf{Attorney} ' s Office, which has used its discretion to drop charges against the two individuals,"Northwestern \textbf{University} officials said Wednesday.
In a statement, the Cook County state ' s \textbf{attorney} ' s office said that in misdemeanor cases, \textbf{police} departments can bypass prosecutors and directly file charges."In @ @ @ @ @ @ @ @ @ @ Office... does not review or approve the charges prior to filing in any instance,"the office said.
State ' s \textbf{Attorney} Kim Foxx had declined to address the situation earlier this week.
The men, 20 and 22, were charged in December with theft of advertising services, a \textbf{class} A misdemeanor, according to Cook County court records.
They were cited, but the charges were still pending as of Tuesday.
Advocates pointed out that criminal charges can follow people for the rest of their lives."These are \textbf{college} \textbf{students} that were engaging in a political \textbf{protest}, one might even describe it as a stunt to make a political point,"Ed Yohnka, spokesman for the American Civil Liberties Union Illinois, told the Tribune Monday."No one was meaningfully harmed as a result of this."Rifqa Falaneh, a fellow at Chicago-based advocacy group Palestine Legal, said the charges were unprecedented coming out of an institution that values free \textbf{speech}."It causes a chilling impact @ @ @ @ @ @ @ @ @ @ believe in,"Falaneh said."There has been enormous community support -- without the pressure, the charges may not have been dropped."Falaneh, who is helping the Northwestern \textbf{students} responsible for the parody newspaper page with public-facing advocacy, said Palestine Legal has seen an uptick in similar cases."Since October 7, we ' ve received more than 1300 reports of repression and censoring targeting those voicing support for Palestinians in Gaza,"she said."And nearly half of these intakes are \textbf{campus} related -- it ' s been an alarming amount."The state ' s \textbf{attorney} ' s office said it decided to dismiss the charges after a thorough review of the circumstances and engaged in discussions with the complainant, the \textbf{campus} newspaper publisher."Our criminal justice system should only be utilized when there is no other recourse for accountability,"the office said in a statement."Northwestern \textbf{University} and \textbf{campus} \textbf{police} are fully equipped to hold the involved individuals accountable, ensuring that such matters are handled @ @ @ @ @ @ @ @ @ @ context and respectful of \textbf{students} ' rights."The charges stemmed from a parody front page of the Daily Northwestern, which \textbf{Students} Publishing Co. said was attached to real editions of the paper and placed around \textbf{campus}.
On Oct. 25, \textbf{students} on \textbf{campus} could find a single-page flyer that looked similar to the popular student-run newspaper with the headline"Northwestern complicit in genocide of Palestinians"printed across its lower third.
A court filing accused the two men of attaching"an unauthorized replica of the Daily Northwestern Newspaper"to a previously distributed edition and placing copies in the newspaper stand.
The charges say they did so"without a contractual agreement between the publisher and an advertiser."In a statement on Monday, \textbf{Students} Publishing Company, parent company of The Daily Northwestern, said it reported the fake front page to \textbf{campus} \textbf{police}, which resulted in charges filed by the Cook County state ' s \textbf{attorney} ' s office.
The next day, it released another statement that said it hired legal @ @ @ @ @ @ @ @ @ @ County State ' s \textbf{Attorney} ' s Office to pursue a resolution to this matter that results in nothing punitive or permanent."In that statement, the publishing company elaborated on the events that led up to the charges, saying it contacted Northwestern \textbf{police} with the intent of protecting \textbf{student} journalists."This co-opting of the work of our \textbf{student} journalists and the potential damage to the reputation of the paper built upon more than a century of hard work was the problem,"the statement said."To us, it seemed no different from someone hacking into our website to post their own content and replace ours."When \textbf{police} identified two people who may have been involved, the publishing company said it signed complaints.
In the statement, the board of directors said they"didn't understand how these complaints started a process that we could no longer control.""We have been listening to our fellow community members, and they have been heard,"the statement read."We @ @ @ @ @ @ @ @ @ @"

% matched lemmas: attorney, campus, class, college, police, protest, school, speech, student, university
\end{textsample}
